% MORSE manuscript draft (2026-10-04). Uses the parent paper's preamble; added references are in morse_new_refs.bib.
% [X] marks results to be filled in from the final experiments.

\begin{abstract}
Cross-embodiment offline reinforcement learning (RL) trains a single control policy on pooled datasets from many
robots, exploiting both expert and abundant suboptimal trajectories. As the proportion of suboptimal data and the
diversity of robots grow, however, the gradients of different robots increasingly conflict and impede learning, most
severely for robots whose morphology differs from the majority. Existing remedies operate on whole robots, while every
joint of every robot is still computed by the same parameters. In this work, we argue that parameter sharing should be
decided at the level of joints, since two robots may benefit from sharing the computation of their knees while their
hips and ankles require different computation. We propose Morphology-aware Routing of Shared Experts (MORSE), which
replaces the whole-body core and the action decoder of a morphology-agnostic policy with pools of experts shared across
robots. A router conditioned on the morphology and state of each joint selects a few experts for that joint in both
pools, so that every joint composes its own whole-body representation and action, and two joints share computation only
through the experts they both select. A mutual-information objective between joints and experts keeps the routing
consistent and all experts in use. On a cross-embodiment offline RL benchmark spanning 16 legged robots, [X].
\end{abstract}

% =====================================================================================================================
\section{Introduction}
\label{sec:intro}
% =====================================================================================================================

Training a single model on large and diverse data has driven much of the recent progress in machine learning, and
robot learning is following the same path. Cross-embodiment learning pools trajectories from many robots to train one
policy that controls all of them, so that experience collected on one platform benefits the
others~\citep{oxe2024,octo2024}. For robots whose bodies differ in the number and arrangement of their joints,
morphology-agnostic policies make this possible by processing every joint with shared modules, conditioned on a
description of the body~\citep{huang2020smp,gupta2022metamorph,bohlinger2024onepolicy}. Offline reinforcement learning
(RL) further broadens the data that such a policy can use, from expert demonstrations to the abundant suboptimal
trajectories that are far cheaper to collect~\citep{levine2020offline,kostrikov2022iql}, and its combination with
cross-embodiment learning has recently been shown to pre-train locomotion policies for many legged
robots~\citep{abe2026crossembodimentoffline}.

Robot bodies, however, are alike in some parts and different in others. The knees of a quadruped and of a humanoid
bend in similar ways and may benefit from shared control, whereas the hip of the quadruped and the ankle of the
humanoid serve different purposes, and even within one robot, the joints of a leg play distinct roles in the gait. A
cross-embodiment policy should therefore share computation between joints that are alike and keep it apart between
joints that differ. Morphology-agnostic policies, in contrast, process every joint of every robot with the same
parameters, so that the update made for one joint changes the computation of all others.

This full sharing becomes a limitation when the data of different robots pull the shared parameters in different
directions. \citet{abe2026crossembodimentoffline} found that, as the proportion of suboptimal data and the number of
robot types grow, the policy gradients of different robots conflict more often, which induces negative transfer and
stalls the learning of robots with few similar embodiments, such as bipeds trained alongside many quadrupeds. Existing
remedies operate on whole robots: gradient surgery projects the conflicting gradients of different
robots~\citep{yu2020pcgrad}, and Embodiment Grouping updates the policy one group of morphologically similar robots at a
time~\citep{abe2026crossembodimentoffline}. Yet the gradient of a robot is a sum of contributions from its joints, so a
conflict between two robots mixes pairs of joints that disagree with pairs that agree, and remedies that treat each
robot as a whole cannot separate them. Specializing computation below the level of whole robots has proven useful
for universal morphology control with limb-specific parameters~\citep{xiong2023modumorph} and for modular robots with
module-wise experts~\citep{liu2026m2oe}, which suggests deciding the sharing for every joint.

To this end, we introduce MORSE (\emph{Morphology-aware Routing of Shared Experts}), a cross-embodiment policy that
learns parameter sharing at the joint level with sparse mixtures of experts
(MoE)~\citep{shazeer2017outrageously,fedus2022switch}. MORSE replaces the whole-body core and the action decoder of a
morphology-agnostic policy~\citep{bohlinger2024onepolicy} with pools of experts shared across robots, and a router that
reads the morphology descriptor and the state of each joint selects a few experts for that joint in each pool. Every
joint thus composes its own representation of the whole body and predicts its action with its own experts. Because the
routing is sparse, the actor gradient of a joint reaches only the experts it selects: joints whose updates conflict can
be served by different experts, while corresponding joints of similar robots can select the same experts and transfer
what they learn (Fig.~\ref{fig:teaser}). A mutual-information objective between joints and
experts~\citep{chen2023modsquad} keeps the routing of every joint consistent across states and all experts in use.
Our main contributions are summarized as follows:
\begin{itemize}
\item We show that the gradient conflicts between
robots in a morphology-agnostic policy decompose into interactions between pairs of joints, and that sparse routing
confines these interactions to the experts that two joints both select.
\item We propose MORSE, in which every joint of every robot
selects shared experts in both the whole-body core and the action decoder according to its morphology and state, so
that every joint learns its own representation of the body and its own path to its action, and parameter sharing is
decided joint by joint rather than robot by robot.
\item We evaluate MORSE on the benchmark
of~\citet{abe2026crossembodimentoffline} with 16 legged robots against robot-level grouping, gradient surgery, and
mixture-of-experts baselines, and analyze which joints share experts across robots.
\end{itemize}

\begin{figure}[t]
\centering
\fbox{\rule{0pt}{1.4in}\rule{0.95\linewidth}{0pt}}
\caption{Learning to share at the joint level. Robot-level remedies update or group whole robots, so all joints share
the same parameters. MORSE lets every joint select experts from pools shared across robots: corresponding joints of
different robots can share experts, while other joints specialize. The selection depends on the morphology and state
of each joint.}
\label{fig:teaser}
\end{figure}

% =====================================================================================================================
\section{Related Work}
\label{sec:related}
% =====================================================================================================================

\subsection{Cross-Embodiment Learning}
Cross-embodiment learning trains a single policy on data from multiple robots so that control knowledge transfers
across platforms. Large-scale efforts pool manipulation data from many robots~\citep{oxe2024,octo2024}, usually with
embodiment-specific input and output modules. For robots whose joints and limbs differ in number and arrangement,
morphology-agnostic policies instead share their modules across joints or limbs, using graph
networks~\citep{wang2018nervenet,huang2020smp}, transformers over limbs~\citep{kurin2021amorpheus,gupta2022metamorph},
or attention over joint descriptors, as in URMA~\citep{bohlinger2024onepolicy}, which controls many legged robots with
one network. We build on URMA because the cross-embodiment offline RL benchmark~\citep{abe2026crossembodimentoffline}
adopts it, which lets us compare directly with the established baselines. In these architectures, every joint of every
robot is computed by the same parameters. To specialize computation below the level of whole robots,
ModuMorph~\citep{xiong2023modumorph} generates the input and output layers of every limb with a hypernetwork
conditioned on the morphology of the limb, and M2oE~\citep{liu2026m2oe} selects experts for each module of a modular
robot with an attention-based gate. In ModuMorph, all limb-specific parameters are produced by one shared hypernetwork,
so the updates of all limbs still meet in its weights, and the generated parameters do not depend on the state; M2oE
addresses the reconfigurations of a modular robot. MORSE targets cross-embodiment offline RL over heterogeneous robots,
where the interference between robots grows with the share of suboptimal
data~\citep{abe2026crossembodimentoffline}. Every joint of every robot selects a few experts from pools shared across
robots, conditioned on its morphology and state, so that two joints share expert parameters only when they select the
same experts, and the selection covers both the whole-body representation and the action decoder.

\subsection{Gradient Conflicts and Parameter Sharing}
When several tasks share a network, their gradients can conflict and cause negative transfer. Multi-task optimization
methods modify the gradients of conflicting tasks, by projecting each onto the normal plane of the
others~\citep{yu2020pcgrad}, by seeking a conflict-averse direction near the average gradient~\citep{liu2021cagrad},
or by bargaining between tasks~\citep{navon2022nashmtl}, and task-grouping methods decide which tasks are updated
together~\citep{standley2020which,fifty2021tag,jeong2025sel}. In cross-embodiment offline RL,
\citet{abe2026crossembodimentoffline} measured conflicts between the actor gradients of different robots and
introduced Embodiment Grouping (EG), which clusters robots by the distance between their morphology graphs and updates
the actor one cluster at a time. These methods keep a single set of parameters shared by all tasks and change only how
or when it is updated. Other methods change which parameters the tasks share: Recon turns the layers with the most
conflicts into task-specific layers~\citep{shi2023recon}, and HarmoDT learns a parameter mask for every task in
multi-task offline RL~\citep{hu2024harmodt}. Both assign parameters to whole tasks identified by their index. MORSE
instead decides the sharing for every joint, learns it end to end from the actor loss without estimating pairwise
gradient conflicts, and conditions it on morphology descriptors rather than task indices, so that corresponding joints
of different robots can share parameters while the joints of one robot can use different ones.

\subsection{Mixtures of Experts}
Sparse mixtures of experts (MoE) route every input to a few experts of a large network and have scaled language and
vision models at a constant cost per
input~\citep{shazeer2017outrageously,lepikhin2021gshard,fedus2022switch,riquelme2021vmoe}. In multi-task learning,
experts shared across tasks are selected by task-dependent gates, so that
related tasks use common experts while unrelated tasks use different ones~\citep{ma2018mmoe,chen2023modsquad}, and in
RL, modules and experts are composed per task for multi-task
manipulation~\citep{yang2020softmodularization,sodhani2021care,sun2022paco,hendawy2024moore} or selected from the
state of a single legged robot across terrains and gaits~\citep{huang2025moeloco}. Experts have also been used to
separate conflicting gradients: STGC detects the tokens of a vision-language model whose gradients conflict within
their expert and reroutes them to other experts~\citep{yang2025stgc}. In these works, the units of routing are tokens,
states, or tasks within a single model or robot. MORSE introduces routing at the level of joints across embodiments.
The joints of different robots select experts according to their morphology and state, so that corresponding joints
of different robots can share experts while the joints of one robot can use different ones, and every joint selects
its own experts in both the whole-body representation and the action decoder. The sharing of parameters thus becomes
a learned decision for every pair of joints, which MORSE obtains end to end from the actor loss, without estimating
gradients.

% =====================================================================================================================
\section{Preliminaries}
\label{sec:prelim}
% =====================================================================================================================

\paragraph{Problem setting.}
We consider a set $\mathcal{R}$ of robots with different morphologies. Robot $r$ has actuated joints $\mathcal{J}_r$
and feet $\mathcal{F}_r$, and every joint $j$ and foot $f$ has a fixed-length descriptor $d_j$ or $d_f$ derived from
the robot's model, such as the position, rotation axis, and actuation limits of a joint; we write $m_r$ for the
descriptors of robot $r$. A state $s$ comprises local observations $o_j$ and $o_f$ of the joints and feet and general
observations $o^{\mathrm{g}}$, such as the velocity of the base and the command, and an action
$a=(a_j)_{j\in\mathcal{J}_r}$ sets a target for every joint. Following \citet{abe2026crossembodimentoffline}, we learn
a single policy $\pi_\theta(a\mid s,m_r)$ from a pooled offline dataset
$\mathcal{D}=\bigcup_{r\in\mathcal{R}}\mathcal{D}_r$ of mixed quality, with every robot contributing equally to the
actor loss:
\begin{equation}
\mathcal{L}(\theta)=\sum_{r\in\mathcal{R}}\mathcal{L}_r(\theta),\qquad
\mathcal{L}_r(\theta)=\frac{1}{|\mathcal{R}|\,n}\sum_{i\in\mathcal{B}_r}\ell_i(\theta),
\label{eq:actor_loss}
\end{equation}
where $\mathcal{B}_r$ holds $n$ transitions of robot $r$ and $\ell_i$ is the per-sample actor loss of the base
learner, such as the advantage-weighted log-likelihood of IQL~\citep{kostrikov2022iql}.

\paragraph{Morphology-agnostic policy.}
URMA~\citep{bohlinger2024onepolicy} encodes the local observation of every joint and foot with its descriptor and
aggregates them, together with $o^{\mathrm{g}}$, into a vector $x$ of fixed size. A whole-body core $C$ maps $x$ to a
representation $z$, which a decoder $D$, shared by all joints, turns into the action mean of every joint:
\begin{equation}
x=A(s,m_r),\qquad z=C(x),\qquad \mu_j=D(z,c_j),
\label{eq:urma}
\end{equation}
where $A$ denotes the encoders and the aggregation, and $c_j$ holds the local encoding $z^{\mathrm{loc}}_j$ of joint
$j$ and an embedding of its descriptor. Every joint of every robot thus reads the same representation $z$ and is
decoded by the same network $D$.

\paragraph{Gradient conflicts between robots.}
Let $\phi$ denote the parameters of $C$ and $D$. Their gradient for robot $r$ is a sum of contributions from its
joints, $g_r=\nabla_\phi\mathcal{L}_r=\sum_{u\in\mathcal{U}_r}g_u$, where $\mathcal{U}_r=\mathcal{B}_r\times
\mathcal{J}_r$ pairs every transition of robot $r$ with every joint and $g_u=(\partial\ell_i/\partial\mu_j)\,
\nabla_\phi\mu_j$ for $u=(i,j)$. The inner product with which \citet{abe2026crossembodimentoffline} measure the
conflict between two robots therefore decomposes over pairs of joints,
\begin{equation}
\bigl\langle g_r,g_{r'}\bigr\rangle=\sum_{u\in\mathcal{U}_r}\sum_{v\in\mathcal{U}_{r'}}\bigl\langle g_u,g_v\bigr\rangle,
\label{eq:conflict}
\end{equation}
so a conflict between two robots mixes pairs of joints that disagree with pairs that agree, and remedies that project,
reweight, or regroup the gradients of whole robots act on these sums as a whole.

% =====================================================================================================================
\section{Morphology-aware Routing of Shared Experts}
\label{sec:method}
% =====================================================================================================================

MORSE turns the two shared networks of Eq.~\eqref{eq:urma}, the whole-body core $C$ and the action decoder $D$, into
pools of $K$ experts that are shared across robots, and every joint selects $k$ experts from each pool. The selection
is made by routers that read the morphology and the state of the joint, so that joints with similar roles can select
the same experts while other joints use different ones, and a mutual-information objective keeps the selection of
every joint consistent and all experts in use. The encoders, the aggregation, and the prediction of the action
standard deviation remain as in URMA, so that MORSE is a drop-in replacement for the core and the decoder.
Figure~\ref{fig:overview} gives an overview.

\begin{figure*}[t]
\centering
\fbox{\rule{0pt}{1.6in}\rule{0.95\textwidth}{0pt}}
\caption{Overview of MORSE. The encoders and the aggregation of a morphology-agnostic policy produce the whole-body
input $x$, which every core expert processes once per state. Every joint mixes the outputs of the $k$ core experts
selected by its router into its own body representation $z_j$, and predicts its action mean with the $k$ decoder
experts selected by a second router. Both routers read the morphology and the state of the joint, and all experts and
routers are shared across robots.}
\label{fig:overview}
\end{figure*}

\subsection{Joint-Level Experts}
\label{sec:experts}

A sparse mixture-of-experts layer outputs a weighted sum of $K$ experts, whose weights come from a router that keeps
only the $k$ largest weights and sets the others to zero~\citep{shazeer2017outrageously}. MORSE applies such layers to
every joint separately. In URMA, every joint reads the same whole-body representation $z$, although the features
of the body that matter to a hip differ from those that matter to an ankle. In MORSE, every joint $j$ instead mixes
$K$ core experts $C_e$, which all read the whole-body input $x$, with the weights $G_{\mathrm{C}}(j)$ given by the
core router:
\begin{equation}
z_j=\sum_{e=1}^{K}G^e_{\mathrm{C}}(j)\,C_e(x).
\label{eq:core}
\end{equation}
The experts depend only on the state, so each is evaluated once per state and reused by every joint that selects it,
while $z_j$ is a view of the whole body that is specific to the joint. The action mean of the joint is produced by
$K$ decoder experts $D_e$, mixed with the weights $G_{\mathrm{D}}(j)$ of the decoder router:
\begin{equation}
\mu_j=\sum_{e=1}^{K}G^e_{\mathrm{D}}(j)\,D_e(z_j,c_j),
\label{eq:decoder}
\end{equation}
with $c_j$ as in Eq.~\eqref{eq:urma}. Separate routers for the two pools let joints share the body features they read
while acting through different decoders, and vice versa. Together, Eqs.~\eqref{eq:core} and~\eqref{eq:decoder} give
every joint its own sparse path from the state to its action.

\subsection{Routing from Morphology and State}
\label{sec:routing}

The router of pool $b\in\{\mathrm{C},\mathrm{D}\}$ is shared by all joints of all robots. It scores the experts from
the morphology and the state of the joint and keeps the $k$ most probable ones:
\begin{equation}
\begin{aligned}
h_b(j)&=\psi^{\mathrm{m}}_b\bigl(d_j,\bar d_r\bigr)+\psi^{\mathrm{s}}_b\bigl(s^b_j\bigr),\\
G_b(j)&=\operatorname{TopK}\bigl(\operatorname{Softmax}(h_b(j))\bigr),
\end{aligned}
\label{eq:router}
\end{equation}
where TopK keeps the $k$ largest entries, sets the others to zero, and renormalizes the kept entries to sum to one.
The morphology branch $\psi^{\mathrm{m}}_b$ reads the descriptor $d_j$ of the joint together with a descriptor
$\bar d_r$ of its robot, formed from the mean joint and foot descriptors and the numbers of joints and feet. Because
all descriptors share one format, corresponding joints of different robots, such as the knees of two quadrupeds,
receive similar descriptors, and this branch gives them similar preferences over experts. The state branch
$\psi^{\mathrm{s}}_b$ reads $s^{\mathrm{C}}_j=[x,o_j]$ in the core and $s^{\mathrm{D}}_j=[z_j,z^{\mathrm{loc}}_j]$ in
the decoder, which lets a joint adjust its selection to its situation, for example between stance and swing. Both
branches are two-layer perceptrons whose output layers start from small random weights, so that different joints
begin with different expert sets. Because the routing is sparse, the actor gradient of a joint reaches only the experts
it selects, so that joints whose updates conflict can be served by different experts, while corresponding joints of
similar robots can share experts and transfer what they learn (Sec.~\ref{sec:analysis}).

\subsection{Mutual Information between Joints and Experts}
\label{sec:mi}

For experts to specialize, every joint should rely on a consistent set of experts across states, and all experts
should be used. We encode both requirements in the mutual information between joints and
experts~\citep{chen2023modsquad}. Let $t=(r,j)$ index the joints of all robots, $N=\sum_{r\in\mathcal{R}}
|\mathcal{J}_r|$ be their number, and $p_{b,i}(j)=\operatorname{Softmax}(h_b(j))$ be the routing probabilities of
joint $j$ in transition $i$ before the top-$k$ selection. For pool $b$, we estimate the probability that joint $t$
uses expert $e$ from the transitions $\mathcal{B}_r$ of its robot in a minibatch, and weight all joints equally:
\begin{equation}
P_b(e\mid t)=\frac{1}{|\mathcal{B}_r|}\sum_{i\in\mathcal{B}_r}p^e_{b,i}(j),\qquad P(t)=\frac{1}{N},
\label{eq:routing_marginals}
\end{equation}
with the marginal $P_b(e)=\sum_tP(t)\,P_b(e\mid t)$. The mutual information between the joint $J$ and the expert $E$,
\begin{equation}
I_b(J;E)=\sum_{t,e}P(t)\,P_b(e\mid t)\log\frac{P_b(e\mid t)}{P_b(e)},
\label{eq:mi}
\end{equation}
is high when every joint concentrates on a few experts across states while all experts are used. Because it uses the
probabilities of all $K$ experts, it also guides the scores of the experts that a joint does not currently select. We
evaluate the routers on detached inputs for this term, so that it shapes the routing while the representations are
learned from the actor loss.

\subsection{Training and Inference}
\label{sec:training}

MORSE is trained by adding the mutual information of both pools to the actor loss of the base learner,
\begin{equation}
\mathcal{L}_{\mathrm{train}}=\mathcal{L}-\lambda\bigl(I_{\mathrm{C}}(J;E)+I_{\mathrm{D}}(J;E)\bigr),
\label{eq:train}
\end{equation}
where $\mathcal{L}$ is the actor loss of Eq.~\eqref{eq:actor_loss} and $\lambda$ weighs the regularizer. The
mutual information is computed on all transitions of every actor minibatch, and the critics keep the architecture
and the losses of the base learner. At inference, the core experts are evaluated once per state, every joint mixes
the outputs of its $k$ selected core experts and of its $k$ selected decoder experts, and the regularizer is used
only during training.

% =====================================================================================================================
\section{Experiments}
\label{sec:experiments}
% =====================================================================================================================

\subsection{Analysis of Joint Interactions}
\label{sec:analysis}

Joint-level routing changes how the updates of different joints interact. Since a joint updates only the experts it
selects, the interaction between two joints $u$ and $v$ in Eq.~\eqref{eq:conflict}, restricted to the experts of a
pool, is a sum over the experts they both select:
\begin{equation}
\bigl\langle g_u,g_v\bigr\rangle=\sum_{e\in\mathcal{S}_u\cap\mathcal{S}_v}G^e(u)\,G^e(v)\,
\bigl\langle g^e_u,g^e_v\bigr\rangle,
\label{eq:localized}
\end{equation}
where $\mathcal{S}_u$ is the set of experts selected by joint $u$, $G^e(u)$ its routing weight, and $g^e_u$ the
gradient of its loss with respect to expert $e$ at unit weight. Joints that select different experts therefore do not
interact through the experts, whereas a shared network, the case $K=k=1$, couples every pair of joints. [X: measured
conflicts between robots on the shared core and decoder of URMA and on the experts of MORSE, and the experts that
joints share across robots.]
